\section*{Chapter \ref{ch:g2:pushes-and-pulls} --- Pushes and Pulls: Forces}

\begin{solution}{exo:g2:pushes-and-pulls:1}
Push; pull; push; pull; push (a squeeze is pushing from both sides).
\end{solution}

\begin{solution}{exo:g2:pushes-and-pulls:2}
Where the \omterm{def:g2:pushes-and-pulls:force}{force} grabs the object, which way it acts, and how strong
it is (by the arrow's length).
\end{solution}

\begin{solution}{exo:g2:pushes-and-pulls:3}
Start, stop, turn, change shape. Examples vary --- for instance:
flicking a marble (start), catching a frisbee (stop), bouncing a
ball off a wall (turn), squashing a snowball (shape).
\end{solution}

\begin{solution}{exo:g2:pushes-and-pulls:4}
Stop it: push against its motion, finger toward the marble. Speed it
up: push the way it already rolls, from behind. Swerve it: push
sideways, across its path.
\end{solution}

\begin{solution}{exo:g2:pushes-and-pulls:5}
Spring back: the sponge and the rubber band. Keep the dent:
plasticine and soft snow.
\end{solution}

\begin{solution}{exo:g2:pushes-and-pulls:6}
The gentle push: a short arrow pointing right. The hard push: a long
arrow pointing left. Different lengths, opposite directions.
\end{solution}

\begin{solution}{exo:g2:pushes-and-pulls:7}
The Earth's pull. It does not need to touch --- it pulls the spoon
through the \omterm{def:g2:air-around-us:air}{air}, and always downward.
\end{solution}

\begin{solution}{exo:g2:pushes-and-pulls:8}
The Earth's pull (a dropped sock falls) and the \omterm{def:g1:magnets:magnet}{magnet}'s pull (the
refrigerator \omterm{def:g1:magnets:magnet}{magnet} holds a drawing through the paper).
\end{solution}

\begin{solution}{exo:g2:pushes-and-pulls:9}
Two arrows on the rope, equally long, pointing opposite ways. The
ribbon stands still when the two pulls are \emph{equally strong} ---
they cancel each other out. The moment one team pulls harder, its
arrow wins and the ribbon moves that way.
\end{solution}

\begin{solution}{exo:g2:pushes-and-pulls:10}
The same \omterm{def:g2:pushes-and-pulls:force}{force} all three times: the Earth's pull, arrow pointing
straight down --- on the way up (slowing the ball), at the top
(about to bring it back), and on the way down (speeding it up). It
never stops pulling for a single moment.
\end{solution}
